\documentclass[12pt]{article}
\usepackage[utf8]{inputenc}
\usepackage[T1]{fontenc}
\usepackage{hyperref}
\title{Ngalum Language Resources for K-12 Curriculum Development}
\author{Compiled by Claude Sonnet 5.5 (claude-sonnet-5-5) for Andr\'as Kornai}
\date{}
\begin{document}
\maketitle

\section{Introduction}
Ngalum (ISO 639-3 szb; also called Sibil) is the most populous of the Ok languages, spoken in the Ok Sibil, Ok Tsop and Ok Nangul valleys of Highland Papua, Indonesia, and across the border in Papua New Guinea. Healey (1964) described three varieties: Sibil, Ngalum proper and Kupel. Kupel was later reclassified as a Ketengban (Mek) variety. Web search found very little material that is usable in teaching. No school materials, primers or textbooks turned up.

\section{Primary Educational Resources}
No primer, textbook or school programme for Ngalum was found in search results. A K-12 curriculum would have to be built from scratch.

\section{Linguistic and Reference Materials}
The New Guinea World survey lists Hylkema's grammar and dictionary of Ngalum (dated 1996) and an unpublished Kroneman vocabulary as the main sources. It also cites Healey (1964), which has comparative vocabulary for the Sibil variety. Another article describes Hylkema's dictionary and grammar as unpublished. I found no digitized or purchasable copy of either. Overview pages: \url{https://newguineaworld.linguistik.uzh.ch/families/trans-new-guinea/central-west-new-guinea/digul-river-ok/ok/ngalum}, \url{https://en.wikipedia.org/wiki/Ngalum_language}.

\section{Religious and Community Resources}
No Ngalum scripture translation turned up in search. For the community, see the Wikipedia page \url{https://en.wikipedia.org/wiki/Ngalum_people} and an Indonesian-language article on Ngalum Ok traditions in Pegunungan Bintang: \url{https://www.goodnewsfromindonesia.id/2025/03/06/suku-ngalum-ok-penjaga-tradisi-leluhur-di-pegunungan-bintang}.

\section{Digital and Audio Resources}
No audio, video or corpus resources were found. The Endangered Languages Project has a catalogue entry: \url{https://endangeredlanguages.com/lang/10696?hl=pt-br}.

\section{Curriculum Development Context}
Little exists for Ngalum. Any curriculum would depend on documentation work with the community, probably starting from the Hylkema materials if they can be located. Teaching in Ngalum is not known to happen in schools.
\end{document}